\documentclass[20pt,margin=5mm, innermargin=3mm]{tikzposter}
\usepackage[utf8]{inputenc}
\geometry{paperwidth=42.0cm,paperheight=59.4cm}
% \usetheme{Simple}
% \usecolorstyle{Spain}
% \usebackgroundstyle{Empty}
% \usetitlestyle{Filled}
% \useblockstyle{Basic}
\tikzposterlatexaffectionproofoff
\usepackage{amsmath}

\useblockstyle{Slide}
\renewcommand*\familydefault{cabin}
\usepackage[sfdefault,condensed]{cabin}
\usepackage[T1]{fontenc}
\usepackage{fix-cm}
\usepackage{tikz,multicol}
\usepackage{capt-of}

\definecolor{greenone}{RGB}{21,75,52}
\definecolor{greentwo}{RGB}{25,97,66}
\definecolor{greenthree}{RGB}{16,57,39}
\definecolor{mygray}{RGB}{203,203,203}
\definecolorstyle{myColorStyle} {
    \colorlet{colorOne}{greenone}
    \colorlet{colorTwo}{greentwo}
    \colorlet{colorThree}{greenthree}
    \colorlet{colorFour}{mygray}
    }{
    % Background Colors
    \colorlet{backgroundcolor}{colorOne}
    \colorlet{framecolor}{colorOne}
    % Title Colors
    \colorlet{titlefgcolor}{white}
    \colorlet{titlebgcolor}{colorTwo}
    % Block Colors
    \colorlet{blocktitlebgcolor}{colorFour}
    \colorlet{blocktitlefgcolor}{colorTwo}
    \colorlet{blockbodybgcolor}{white}
    \colorlet{blockbodyfgcolor}{black}
    % Innerblock Colors
    \colorlet{innerblocktitlebgcolor}{white}
    \colorlet{innerblocktitlefgcolor}{black}
    \colorlet{innerblockbodybgcolor}{white}
    \colorlet{innerblockbodyfgcolor}{black}
    % Note colors
    \colorlet{notefgcolor}{black}
    \colorlet{notebgcolor}{white}
    \colorlet{notefrcolor}{white}
}

\usecolorstyle{myColorStyle}

\usebackgroundstyle{Rays}
\usepackage{adjustbox}
\usetitlestyle{Empty}

\makeatletter %To make columns scale to custom width 
\setlength{\TP@visibletextwidth}{\textwidth-2\TP@innermargin}
\setlength{\TP@visibletextheight}{\textheight-2\TP@innermargin}
\makeatother

 
%\graphicspath{{images/}}

%%%%%%%%%%%%%%%%%%%%%%%%%%%%%%%%%%%%%%%%%%%%%%%%%%%%%%%%%%%%%%%%%%%%%%%%

\title{\parbox{\linewidth}{\vspace{-4.5cm} \centering \huge Scaling up individual metabolism to ecosystem fluxes}}
\author{\vspace{-100pt} % space between title and author
  {\begin{minipage}{8cm}
    \hfill
    \vspace{-20pt}
    \flushleft
    \includegraphics[trim=0cm 8cm 0cm 8cm,clip,width=5cm]{Imperial_Whitelogo.pdf}
  \end{minipage}
    \vspace{-20pt}
    \hspace{2pt}}
  {\begin{minipage}{20cm}
    \centering
        { \Large Sofia Sal \& Samraat Pawar \\ \vspace{-20pt}
          s.sal-bregua@imperial.ac.uk  }
  \end{minipage}
    \vspace{-20pt}
  \hspace{2pt}}
  {\begin{minipage}{8cm}
    \centering
    \hfill
    \centering
    \includegraphics[trim=0cm 8cm 0cm 8cm,clip,width=8cm]{ExeterLogo.pdf}
  \end{minipage}
  \vspace{-20pt}}
  }

%%%%%%%%%%%%%%%%%%%%%%%%%%%%%%%%%%%%%%%%%%%%%%%%%%%%%%%%%%%%%%%%%%%%%%%%
\makeatletter
\newcommand\semiHUGE{\@setfontsize\semiHUGE{60}{27.38}}
\makeatother

\begin{document}
\node[xshift=-20cm, at=(topright),opacity=0.9]{\includegraphics[width=\paperwidth, height=10cm]{head1low.jpg}} ;

\maketitle[titletoblockverticalspace=-8mm]

  %%%%%%%%%%%%%%%%%%%%%%%%%%%%%%%%%%%%%%%%%%%%%%%%%%%%%%%%%%%%%%%%%%%%%%%%

\block[titleinnersep=5mm]{\small Do species-level thermal performance curves (TPCs) scale up to ecosystem fluxes?}{

  \begin{minipage}[]{0.45\linewidth}
    \centering
    \includegraphics[trim=6cm 2cm 4cm 6cm,clip,width=5cm]{EcoModelFigure2.pdf}
    
  \end{minipage}   
  \hspace{.8cm}    
  \begin{minipage}[]{0.52\linewidth}

    %% {\bf Scaling up} is a linear mapping (strictly what geometric scaling means) between the lower-level phenomenon (here,individual-level or intra-specific thermal responses) and the emergent one of ecosystem functioning. \\

   Previous studies assume that the temperature dependence of ecosystem function is a simple scaling up of all the component species’ thermal responses. In this case, predicting the effects of climatic warming or cooling on ecosystem function would be a relatively straightforward task. 

   
    \end{minipage}
 }

\begin{columns}
  %%%%%%%%%%%%%%%%%%%%%%%%%%%%%%%%%%%%%%%%%%%%%%%%%%%%%%%%%%%%%%%%%%%%%%%%
  \column{0.45}
  \block{Parametrization}{
    
    \centering
    \includegraphics[width=\linewidth]{IntraspPlots.pdf}
    \footnotesize Representative set of individual-level plot pairs for R (left) and NP (right). Each pair of R-NP corresponds to the same experiment and species. These species are associated to temperate regions.\\
    
     \vspace{20pt}
    \includegraphics[width=\linewidth]{Terr_EaPairHistInt.pdf}
    \footnotesize Distributions of intraspecific activation energies for R and NP.

  }

  \column{0.55}
  \block{Validation}{ 
    \centering
    \includegraphics[width=\linewidth]{EcosystemPlots.pdf}
    \footnotesize Representative set of ecosystem flux responses to temperature for different sites from tropical and temperate regions. Y axis show respiration in log scale. 

    \includegraphics[width=\linewidth]{Terr_EaPairHistEco.pdf}
    \captionof{figure}{\footnotesize Distributions of activation energies for Reco and NEP at ecosystem level.}
    
  }
\end{columns}


%%%%%%%%%%%%%%%%%%%%%%%%%%%%%%%%%%%%%%%%%%%%%%%%%%%%%%%%%%%%%%%%%%%%%%%%
\begin{columns}

  \column{0.55}
  \block[bodyverticalshift=-9mm]{Summary}{ 	
    \begin{itemize}
      
    \item At the intra-specific level, $E_a$ and $T_{peak}$ for $R$ are usually higher than for $P$.
    \item $T_{peak}$’s are usually much higher than the “characteristic” adaptive environment of the organism, so the full unimodal thermal response doesn't matter for mapping individual TPCs to ecosystem-level fluxes.
     \item Ecosystem responses strongly depend on the temperature fluctuations of each ecosystem regime.
    \end{itemize}         
  }
  %%%%%%%%%%%%%%%%%%%%%%%%%%%%%%%%%%%%%%%%%%%%%%%%%%%%%%%%%%%%%%%%%%%%%%%%
  \column{0.45}
  \block{Ongoing}{
    
    \begin{itemize}\itemsep10pt
    \item What role might acclimation of intraspecific TPCs play on the ecosystem response?
    \item What is an appropriate distribution for species-level biomass abundances $x_i$?
    \item  Which is the effect of non-linear interactions between species? 
    \end{itemize}      
  }
  
\end{columns}

\block[bodyverticalshift=-8mm]{}{
  \footnotesize   
  \begin{itemize}
  \item Enquist, B. J., Economo, E. P., Huxman, T. E., Allen, A. P., Ignace, D. D. and Gillooly, J. F. 2003. Scaling metabolism from organisms to ecosystems. – Nature 423(6940): 639–42.
  \item Savage, V. M. 2004. Improved approximations to scaling relationships for species, populations, and ecosystems across latitudinal and elevational gradients. – J. Theor. Biol. 227(4): 525–534.
  \item Yvon-Durocher, G. and Allen, A. P. 2012. Linking community size structure and ecosystem functioning using metabolic theory. – Philos. Trans. R. Soc. Lond. B. Biol. Sci. 367(1605): 2998–3007.
  \item Yvon-Durocher, et al. 2012. Reconciling the temperature dependence of respiration across timescales and ecosystem types. – Nature 487(7408): 472–6.
  \end{itemize}
  \textit{Acknowledge: This work used eddy covariance data acquired by the FLUXNET community (http://fluxnet.fluxdata.org).}
}


\end{document}
